% TOP_PROBLEM: {"id": 20001047, "title": "One-sided lifting and polarized minimal counterexamples for Sullivan-3", "category_id": 2, "category": {"id": 2, "name": "combinatorics", "display_name": "Combinatorics", "description": "Counting problems, graph theory, discrete structures.", "slug": "combinatorics", "order_index": 2, "created_at": "2026-07-31T15:26:25.671Z"}, "status": "partially_solved", "problem_number": "AIM-COMBINATORICS-0172", "set_id": 14}
% FIRST_POSTED: 2026-10-01
% ENTRY_KIND: statement_only
% UnsolvedMath ID 20001047; source catalogue code AIM-COMBINATORICS-0172
% !TeX program = lualatex
% Definition and sources only; this file is not an LLM solution attempt.
\documentclass[11pt,a4paper]{article}
\usepackage[margin=25mm]{geometry}
\usepackage{fontspec}
\setmainfont{lmroman10-regular.otf}[BoldFont=lmroman10-bold.otf,
 ItalicFont=lmroman10-italic.otf,BoldItalicFont=lmroman10-bolditalic.otf]
\usepackage{amsmath,amssymb,mathtools}
\usepackage{unicode-math}
\setmathfont{latinmodern-math.otf}
\AtBeginDocument{\let\setminus\smallsetminus}
\usepackage{xurl,hyperref}
\hypersetup{colorlinks=true,urlcolor=blue}
\setcounter{secnumdepth}{0}
\setlength{\parindent}{0pt}
\setlength{\parskip}{5pt}
\setlength{\emergencystretch}{3em}
\begin{document}
\section*{One-sided lifting and polarized minimal counterexamples for Sullivan-3}\label{20001047}
{\small UnsolvedMath ID 20001047 · AIM-COMBINATORICS-0172 · Combinatorics · AIM Workshop Problem Lists}
\subsection{Definitions and mathematical statement}
Let $D$ be any nonempty finite simple directed graph without loops
or digons, where a digon is a pair $u\to v$, $v\to u$.
Write $d^+(v)$ and $d^-(v)$ for the outdegree and indegree of a
vertex $v$. Let $N_2^+(v)$ be the set of vertices $w$ for which
there is a directed path $v\to u\to w$, but no arc $v\to w$,
with $w\ne v$. Equivalently, the shortest directed distance from
$v$ to $w$ is exactly two. Several paths with endpoint $w$ contribute
only one member to this set.

The question is whether there must exist a vertex $v$ satisfying
\[
           |N_2^+(v)|+d^+(v)\ge2\min\{d^-(v),d^+(v)\}.
\]
If $d^-(v)\ge d^+(v)$, this is the inequality
$|N_2^+(v)|\ge d^+(v)$. If $d^-(v)\le d^+(v)$, it is instead
$|N_2^+(v)|\ge2d^-(v)-d^+(v)$. The graph need only have one
vertex satisfying the appropriate inequality for its own degrees.
These two cases explain the role of the minimum without imposing
a global relation between in- and outdegrees.

No assumptions on the underlying undirected graph, beyond simplicity,
are made. A source or sink already satisfies the desired inequality.
The statement is Conjecture 6.7 of the source, independently of its
motivating implications from stronger rainbow or second-neighbourhood
conjectures. Its target is this particular mixed-degree bound.
\subsection{Short English statement}
Is there always a vertex whose first two outgoing layers have size at least twice the smaller of its indegree and outdegree?
\subsection{Sources}
\begin{itemize}
\item[S1] ULAM.AI and the UnsolvedMath Contributors, supplied problems.json, record 20001047 (AIM-COMBINATORICS-0172), AIM Workshop Problem Lists. Catalogue identity and source transcription; CC BY 4.0.
\url{https://huggingface.co/datasets/ulamai/UnsolvedMath}.
\item[S2] American Institute of Mathematics, The Caccetta-Haggkvist conjecture, item 6.7. Original workshop question underlying AIM-COMBINATORICS-0172.
\url{https://aimath.org/WWN/caccetta/caccetta.pdf}.
\end{itemize}
\end{document}
